% multifat method sheets: one section per registered criterion, for the FABER WG4.3 review.
% Build: latexmk -pdf -interaction=nonstopmode -halt-on-error methods.tex
\documentclass[10pt,a4paper]{article}
\usepackage[margin=22mm]{geometry}
\usepackage[T1]{fontenc}
\usepackage{lmodern,microtype,amsmath,amssymb,booktabs,array,longtable,xcolor,enumitem}
\usepackage[colorlinks=true,linkcolor=black,citecolor=black,urlcolor=blue!55!black]{hyperref}
\setlength{\parindent}{0pt}
\setlength{\parskip}{4pt}
\setlist{leftmargin=5mm,topsep=2pt,itemsep=1pt}
\numberwithin{equation}{section}

% ---- shared notation (defined in notation.tex, Section 1) ----
\newcommand{\Tsig}{\boldsymbol\sigma}            % Cauchy stress tensor
\newcommand{\Teps}{\boldsymbol\varepsilon}        % elastic strain tensor
\newcommand{\Sm}{\boldsymbol\sigma_m}             % mean tensor
\newcommand{\Ss}{\boldsymbol\sigma_s}             % sine coefficient tensor
\newcommand{\Sc}{\boldsymbol\sigma_c}             % cosine coefficient tensor
\newcommand{\dev}{\operatorname{dev}}
\newcommand{\tr}{\operatorname{tr}}
\newcommand{\n}{\mathbf n}
\newcommand{\MPa}{\,\mathrm{MPa}}
\newcommand{\fm}{f_{-1}}                          % fully reversed axial limit
\newcommand{\tm}{t_{-1}}                          % fully reversed torsion limit

% one block per method; \methodhead{registry name}{title}
\newcommand{\methodhead}[2]{\section{#2 \texorpdfstring{\normalfont\small(\texttt{#1})}{(#1)}}}
\newcommand{\block}[1]{\par\smallskip\textbf{#1.}\ }

\title{Multiaxial fatigue criteria in \texttt{multifat}\\[2pt]\large Method sheets}
\author{FABER WG4.3}
\date{\today}

\begin{document}
\maketitle

\noindent Each sheet states a criterion as its original paper defines it, every symbol it uses, how
its constants are obtained, how a life follows, the loading it covers, and every choice the
repository makes where the paper is silent or departs. The worked example of each sheet is the
one the repository's tests execute. Section~\ref{sec:notation} defines the notation shared by all
sheets; a sheet restates only what is particular to its method.

\tableofcontents

\section{Notation and common procedures}\label{sec:notation}

\subsection{Frame, loading and stress quantities}

\block{Frame} A right-handed orthonormal frame at the critical point of a plain specimen:
$x$ is the circumferential (hoop) direction, $y$ the specimen axis, $z$ the outward surface
normal. A thin-walled or solid specimen under axial force, torsion and internal pressure carries
the plane stress $\sigma_{xx}$ (hoop), $\sigma_{yy}$ (axial), $\tau_{xy}$ (torsion); bending is
read as an axial stress $\sigma_{yy}$ at the surface. Stresses are in MPa, nominal and linear
elastic unless a sheet states otherwise.

\block{Harmonic loading} Every test is a constant-amplitude, single-frequency cycle. Each stress
component is written as
\begin{equation}
\Tsig(t) = \Sm + \Ss\sin\omega t + \Sc\cos\omega t ,\qquad 0\le \omega t < 2\pi ,
\label{eq:harmonic}
\end{equation}
with the symmetric tensors $\Sm$ (mean), $\Ss$ and $\Sc$ (in-phase and quadrature parts). For an
axial amplitude $\sigma_a$ and a torsional amplitude $\tau_a$ with phase shift $\delta$,
$\sigma_{yy}(t) = \sigma_m + \sigma_a\sin\omega t$ and
$\tau_{xy}(t) = \tau_m + \tau_a\sin(\omega t + \delta)$, i.e.\ $\tau_a\cos\delta$ enters $\Ss$ and
$\tau_a\sin\delta$ enters $\Sc$. The loading is \emph{proportional} when $\Ss$ and $\Sc$ are
parallel (as six-component vectors); the mean $\Sm$ may then have any direction.

\block{Scalar of one harmonic} A scalar $q(t) = m + a\sin\omega t + b\cos\omega t$ has
\begin{equation}
q_a = \sqrt{a^2+b^2}\ \text{(amplitude, half range)},\qquad q_m = m,\qquad
q_{\max} = m + q_a ,\qquad q_{\min} = m - q_a .
\label{eq:scalar}
\end{equation}
The stress ratio of a uniaxial cycle is $R = q_{\min}/q_{\max}$; fully reversed means $R=-1$.

\block{Plane quantities} A material plane is identified by its unit normal $\n$ ($\n$ and $-\n$
are the same plane). On it, the normal stress and the shear stress vector are
\begin{equation}
\sigma_n(t) = \n^{\mathsf T}\Tsig(t)\,\n ,\qquad
\boldsymbol\tau(t) = \Tsig(t)\,\n - \sigma_n(t)\,\n .
\label{eq:plane}
\end{equation}
$\sigma_{n,a}$, $\sigma_{n,m}$ and $\sigma_{n,\max}$ follow from Eq.~\eqref{eq:scalar}. The shear
amplitude $\tau_a$ of a non-proportional path (the trace of $\boldsymbol\tau(t)$ is an ellipse)
needs a rule, stated by each critical-plane sheet.

\block{Invariants} $I_1 = \tr\Tsig$ (first invariant), $\sigma_H = I_1/3$ (hydrostatic stress),
$\dev\Tsig = \Tsig - \sigma_H\mathbf I$, $J_2 = \tfrac12\,\dev\Tsig:\dev\Tsig$. For one harmonic
the amplitudes are
\begin{equation}
I_{1,a} = \sqrt{(\tr\Ss)^2 + (\tr\Sc)^2},\qquad
\sqrt{J_{2,a}} = \sqrt{\tfrac12\left(\lVert\dev\Ss\rVert^2 + \lVert\dev\Sc\rVert^2\right)} ,
\label{eq:invariants}
\end{equation}
with $\lVert\cdot\rVert$ the Frobenius norm. For one harmonic, the minimum circumscribed ellipse
(Li), the prismatic hull (Mamiya) and the principal-coordinate prismatic hull (Khalij) of the
deviatoric path all give this $\sqrt{J_{2,a}}$; the minimum circumscribed circle does not.
Uniaxial $\sigma_a$: $\sqrt{J_{2,a}}=\sigma_a/\sqrt3$; torsion $\tau_a$: $\sqrt{J_{2,a}}=\tau_a$.

\block{Elastic strains} Where a criterion needs strains, they are the linear-elastic strains of
the stress, $\Teps = \bigl[(1+\nu)\Tsig - \nu\,\tr(\Tsig)\,\mathbf I\bigr]/E$, unless a sheet states
otherwise. Shear strains $\gamma = 2\varepsilon_{ij}$ are engineering strains.

\subsection{Material constants}

\begin{center}\small
\begin{tabular}{@{}lllp{80mm}@{}}
\toprule
Symbol & Code & Unit & Meaning\\\midrule
$\fm$ & \texttt{Sigm1} & MPa & fully reversed ($R=-1$) axial or bending fatigue limit (amplitude)\\
$\tm$ & \texttt{Taum1} & MPa & fully reversed torsional fatigue limit (amplitude)\\
$f_0$ & \texttt{Sig0} & MPa & repeated ($R=0$) axial limit, given as the \emph{maximum} stress of its cycle\\
$t_0$ & \texttt{Tau0} & MPa & repeated ($R=0$) torsional limit, maximum stress\\
$S_u$, $S_y$ & \texttt{Su}, \texttt{Sy} & MPa & ultimate tensile strength, (monotonic) yield strength\\
$E$, $\nu$ & \texttt{E}, \texttt{nu} & MPa, -- & Young's modulus, Poisson's ratio\\
$\kappa$, $r$ & -- & -- & limit ratios $\kappa = \fm/\tm$ and $r = \tm/\fm = 1/\kappa$\\
$N_\text{ref}$ & -- & cycles & life at which the limits are defined ($10^6$--$10^7$)\\
\bottomrule
\end{tabular}
\end{center}
A sheet lists any further constant with its source; a constant fitted by the repository is
stated as such.

\subsection{From a damage parameter to a life}\label{sec:mastercurve}

\block{Damage parameter} Each criterion maps a loading cycle to a scalar damage parameter
$\mathrm{DP}$ (an equivalent stress, a fatigue index, an energy). A fatigue-limit criterion
returns an index that equals 1 on the limit.

\block{Master curve} Unless the method has its own life law, the life follows from
\begin{equation}
\mathrm{DP} = A\,N_f^{\,b},\qquad N_f = (\mathrm{DP}/A)^{1/b},
\label{eq:mastercurve}
\end{equation}
fitted by least squares in $\log$--$\log$ space on the \emph{design set}: the plain, fully reversed
uniaxial tension failures of the material, with $|R+1|\le 10^{-3}$ and every other stress
component below $10^{-3}$ of the peak (a sheet may state the torsion failures instead). Run-outs
never enter the fit. The curve is used only if it is identifiable: at least 5 tests, at least one
decade of life, and a slope significantly different from zero. A criterion that is a fatigue-limit
condition gains finite lives only through Eq.~\eqref{eq:mastercurve}; this is a repository
extension and is stated as such.

\block{Load-equivalent error} The error of a prediction is measured on the load, so that criteria
with different damage parameters are comparable. For a test of life $N_\text{exp}$, $s^*$ is the
factor on all its stresses (means and amplitudes) at which the model predicts $N_\text{exp}$:
$\mathrm{DP}(s^*\Tsig) = A N_\text{exp}^{\,b}$. Then
\begin{equation}
e_L = \ln(1/s^*) ,
\label{eq:eload}
\end{equation}
$e_L>0$ is conservative (the model reaches its curve at a lower load than the test) and the load is
misjudged by $e^{e_L}-1$. A load beyond a criterion's domain counts as infinite damage; a root that
lies on the edge of that domain is rejected.

\block{Fatigue-limit datasets} For a set of limit states (no failures), the reference is the
damage of the fully reversed axial limit, $\mathrm{DP}_\text{ref}=\mathrm{DP}(\fm)$, and
$\mathrm{FI} = \mathrm{DP}/\mathrm{DP}_\text{ref}-1$. The load-scaled index $\mathrm{FI}_\text{load}=1/s-1$,
with $\mathrm{DP}(s\Tsig)=\mathrm{DP}_\text{ref}$, gives the same $e_L=\ln(1+\mathrm{FI}_\text{load})$.

\subsection{Critical-plane search}\label{sec:planes}

A critical-plane criterion evaluates its damage on candidate planes: a 522-plane grid over the
hemisphere of normals, analytical seed planes, and a local refinement. The selected plane maximises
the criterion's own selection metric; ties go to the larger damage. The repository reports the
damage on that plane and, as the variant \texttt{\_ext}, the maximum damage over all planes.

\subsection{Scope}

A criterion is evaluated only on the loading and material its source covers. Outside it the
repository raises a scope error, the model is reported as skipped, and nothing is extrapolated
with an invented formula.


% ---- method sheets (one file per method in sheets/) ----
% registry: KAKUNO
\methodhead{KAKUNO}{Kakuno and Kawada (1979): octahedral criterion with amplitude and mean terms}\label{sheet:KAKUNO}

\block{Source} H.~Kakuno, Y.~Kawada, \emph{Fatigue Engng Mater.\ Struct.} 2 (1979) 229--236
\cite{kakuno1979}, Eqs.~(4), (6)--(12); Figs.~2--6.

\block{Idea and scope} A correction of Sines' criterion: the octahedral shear stress amplitude at
the fatigue limit decreases linearly with the octahedral normal stress, with \emph{different}
slopes for its amplitude and its mean. Sines keeps only the mean term, which forces
$\tm/\fm=1/\sqrt3$ without mean stresses. Proposed for fatigue limits of round bars under repeated
bending and torsion with superimposed static bending and torsion; plain and moderately notched
members.

\block{Inputs}
\begin{center}\small
\begin{tabular}{@{}llp{100mm}@{}}
\toprule
Symbol & Unit & Meaning\\\midrule
$\sigma_{a_i}$, $\sigma_{m_i}$ ($i=1,2,3$) & MPa & principal stress amplitudes and principal mean stresses\\
$\sigma_a,\sigma_m$; $\tau_a,\tau_m$ & MPa & bending and torsion amplitude and mean (round bar)\\
$\sigma_{wb}=\fm$, $\tau_w=\tm$ & MPa & fully reversed bending and torsion limits\\
$\sigma_{up}=f_0$ & MPa & pulsating ($R=0$) bending limit, maximum stress of the cycle\\
$k_1=\sigma_{up}/(2\sigma_{wb})$ & -- & ratio of the pulsating amplitude to the reversed limit\\
$A$, $B$, $C$ & MPa, --, -- & material constants, Eqs.~\eqref{eq:kk-A}--\eqref{eq:kk-C}\\
\bottomrule
\end{tabular}
\end{center}

\block{Formulation} The criterion (Eq.~4) is
\begin{equation}
\tfrac13\sqrt{(\sigma_{a_1}-\sigma_{a_2})^2+(\sigma_{a_2}-\sigma_{a_3})^2+(\sigma_{a_3}-\sigma_{a_1})^2}
 = A - B\,(\sigma_{a_1}+\sigma_{a_2}+\sigma_{a_3}) - C\,(\sigma_{m_1}+\sigma_{m_2}+\sigma_{m_3}) ,
\label{eq:kk}
\end{equation}
whose left-hand side is the octahedral shear stress amplitude $\tau_{oct,a}=\sqrt{2/3}\sqrt{J_{2,a}}$.
For bending and torsion it becomes $\tfrac13(2\sigma_a^2+6\tau_a^2)^{1/2}=A-B\sigma_a-C\sigma_m$
(Eq.~6). Imposing the reversed torsion limit, the reversed bending limit and the pulsating bending
limit ($\sigma_a=\sigma_m=\sigma_{up}/2$) gives (Eqs.~7--9)
\begin{align}
A &= \sqrt{\tfrac23}\,\tau_w , \label{eq:kk-A}\\
B &= \frac{\sqrt2}{3}\left(\sqrt3\,\frac{\tau_w}{\sigma_{wb}} - 1\right), \label{eq:kk-B}\\
C &= \sqrt{\tfrac23}\,\frac{1-k_1}{k_1}\,\frac{\tau_w}{\sigma_{wb}} . \label{eq:kk-C}
\end{align}
Substituting, the limit curve of combined bending and torsion (Eq.~10) is
\begin{equation}
\Bigl(\frac{\tau_a}{\tau_w}\Bigr)^2\frac{1}{p^2}
+ \Bigl(\frac{\sigma_a}{\sigma_{wb}}\Bigr)^2\frac{\tfrac{2}{\sqrt3}\tfrac{\sigma_{wb}}{\tau_w}-1}{p^2}
+ \frac{\sigma_a}{\sigma_{wb}}\,\frac{2-\tfrac{2}{\sqrt3}\tfrac{\sigma_{wb}}{\tau_w}}{p} = 1 ,
\qquad p = 1 - \frac{\sigma_m}{\sigma_{wb}}\frac{1-k_1}{k_1} ,
\end{equation}
and $p=1$ without mean stresses (Eq.~11). A torsional mean stress does not enter. The paper
recommends $k_1=0.83$ as an average for structural carbon and alloy steels.

\block{Constants and calibration} $A$, $B$, $C$ from $\tau_w$, $\sigma_{wb}$, $\sigma_{up}$; nothing
is fitted.

\block{Life} The paper gives none; finite lives use the master curve on the index below
(Section~\ref{sec:mastercurve}), a repository extension.

\block{Repository choices}
\begin{itemize}
\item The model returns the index
  $\bigl(\tau_{oct,a} + B\,\lvert\tr\boldsymbol\sigma_a\rvert + C\,\tr\Sm\bigr)/A$, equal to 1 on the
  limit, with $\boldsymbol\sigma_a$ the amplitude tensor; the sum of the principal amplitudes is
  taken as $\lvert\tr\boldsymbol\sigma_a\rvert$, positive like any amplitude.
\item Proportional loading only (principal stress amplitudes are undefined otherwise): a case with
  non-parallel $\Ss$, $\Sc$ is out of scope. The mean may have any direction.
\item $\sigma_{wb}$, $\tau_w$, $\sigma_{up}$ are the workbook $\fm$, $\tm$, $f_0$. Without $f_0$, $C$
  is undefined: a case with $\tr\Sm\neq0$ is out of scope; $k_1=0.83$ is not used as a default.
\end{itemize}

\block{Worked example} Ni--Cr--Mo steel of Fig.~5, stresses divided by $\sigma_{wb}$:
$\tau_w=0.617$, $k_1=0.83$, so $A=0.5038$, $B=0.0324$, $C=0.1032$. With $\sigma_m=0.46$,
$p = 1-0.46\cdot0.17/0.83 = 0.906$, and the curve meets the $\tau$ axis at
$\tau_a = p\,\tau_w = 0.559$: $\tau_{oct,a}=\sqrt{2/3}\cdot0.559=0.4563$, index
$(0.4563+0+0.1032\cdot0.46)/0.5038 = 1.000$, with or without a static torsion.

\block{Verification} Eqs.~(7)--(10) and (12) hold to machine precision. Fig.~5: the $\tau$
intercept is 0.559 against 0.557 read from the figure; the $\sigma$ intercept is 0.906 against
about 0.93 drawn (unexplained). The paper prints no table.

% registry: KHALIJ
\methodhead{KHALIJ}{Khalij, Pagnacco and Troian (2017): improved Gough / Nishihara--Kawamoto criterion}\label{sheet:KHALIJ}

\block{Source} L.~Khalij, E.~Pagnacco, R.~Troian, \emph{Int.\ J.\ Fatigue} 99 (2017) 137--150
\cite{khalij2017}, Eqs.~(12), (22), (23), (26); Tables~1 and~6.

\block{Idea and scope} A fatigue-limit criterion for zero-mean sinusoidal plane stress that
weights an equivalent shear stress amplitude against the amplitude of the first invariant. The
weight depends on the limit ratio $r=\tm/\fm$, so one expression moves from a von Mises ellipse
(mild steels, $r\approx 1/\sqrt3$) to a Rankine-like envelope (cast irons, $r\to1$). Proposed for
metals with $0.5\lesssim r\le 1$, in-phase and out-of-phase loading.

\block{Inputs}
\begin{center}\small
\begin{tabular}{@{}llp{100mm}@{}}
\toprule
Symbol & Unit & Meaning\\\midrule
$s_{xx,a}$, $s_{yy,a}$, $s_{xy,a}$ & MPa & amplitudes of the plane-stress components\\
$\phi_{yy}$ & rad & phase shift between the two normal stresses $s_{xx}$ and $s_{yy}$\\
$s_{-1}=\fm$, $\tm$ & MPa & fully reversed axial (or bending) and torsion limits\\
$r=\tm/\fm$ & -- & limit ratio\\
$\alpha(r)$ & -- & weight of the shear term, Eq.~\eqref{eq:khalij-alpha}\\
$J_{2,a}$, $I_{1,a}$ & MPa$^2$, MPa & PHPS shear amplitude squared and first-invariant amplitude\\
$\mathfrak D$ & -- & damage; $\mathfrak D\le1$ means no failure\\
\bottomrule
\end{tabular}
\end{center}

\block{Formulation} With the PHPS amplitude (Eq.~12), which depends on the phase of the normal
stresses only, and the amplitude of the first invariant (Eq.~22),
\begin{gather}
J_{2,a} = \tfrac13\left(s_{xx,a}^2 + s_{yy,a}^2 - s_{xx,a}s_{yy,a}\cos\phi_{yy}\right) + s_{xy,a}^2 ,\qquad
I_{1,a} = \sqrt{s_{xx,a}^2 + s_{yy,a}^2 + 2\,s_{xx,a}s_{yy,a}\cos\phi_{yy}} ,\\
\mathfrak D = \alpha(r)\left[\Bigl(\frac{s_{xx,a}}{\fm}\Bigr)^2 + \Bigl(\frac{s_{yy,a}}{\fm}\Bigr)^2
 - \frac{s_{xx,a}}{\fm}\frac{s_{yy,a}}{\fm}\cos\phi_{yy}\right]
 + \Bigl(\frac{s_{xy,a}}{\tm}\Bigr)^2 + \bigl(1-\alpha(r)\bigr)\frac{I_{1,a}}{\fm} \le 1 ,
 \label{eq:khalij}\\
\alpha(r) = \frac{1}{3r^2} . \label{eq:khalij-alpha}
\end{gather}
Eq.~\eqref{eq:khalij-alpha} is the authors' setting ``OP'': it makes the criterion independent of
the frame for pure shear (Eqs.~24--26). With it the first two groups of Eq.~\eqref{eq:khalij} are
$J_{2,a}/\tm^2$, so
\begin{equation}
\mathfrak D = \frac{J_{2,a}}{\tm^2} + \Bigl(1 - \frac{1}{3r^2}\Bigr)\frac{I_{1,a}}{\fm} .
\label{eq:khalij-invariant}
\end{equation}
Both limits give $\mathfrak D=1$ exactly. For $r<1/\sqrt3$ the weight of $I_{1,a}$ is negative.

\block{Constants and calibration} None fitted: $\alpha$ follows from $\fm$ and $\tm$.

\block{Life} The paper gives none; finite lives use the master curve on $\mathfrak D$
(Section~\ref{sec:mastercurve}), a repository extension. $\mathfrak D$ is not homogeneous in the
load (linear $I_{1,a}$ term).

\block{Repository choices}
\begin{itemize}
\item Eq.~\eqref{eq:khalij-invariant} is evaluated with $J_{2,a}$ and $I_{1,a}$ of
  Eq.~\eqref{eq:invariants}; in plane stress they equal the paper's expressions.
\item Zero mean stress only. The paper's mean-stress extension (Eqs.~29--31) is not implemented:
  its term $\gamma_2(r)\,s_{ms}/s_{-1}^2$ is not dimensionless ($s_{ms}$ is linear in the mean
  stresses), Eq.~(30) for $\gamma_2$ follows from a purely quadratic base criterion, and no reading
  reproduces Table~7. A case with a mean stress, or with an out-of-plane component, is out of scope.
\item $r$ is computed from the material limits; Table~6 uses this unrounded ratio, not the
  two-digit $r$ printed beside it.
\end{itemize}

\block{Worked example} Table~6, CrV-V ($\fm=530$, $\tm=309\MPa$, $r=0.583$, $\alpha=0.981$), biaxial
limit $s_{xx,a}=0.24\fm=127.2$, $s_{yy,a}=-0.84\fm=-445.2\MPa$ (in phase, opposite signs):
$J_{2,a}=\tfrac13(127.2^2+445.2^2+127.2\cdot445.2)=90\,337\MPa^2$, $I_{1,a}=318.0\MPa$,
$\mathfrak D = 90\,337/309^2 + (1-0.981)\,318.0/530 = 0.946+0.012 = 0.958$ (printed 0.95).

\block{Verification} Table~6, OP column: 51 of 53 limits to one unit of the last printed digit;
two 3\%C cast-iron rows (1.035 and 1.063 against 1.06 and 1.01) are unexplained.

% registry: MATSUBARA
\methodhead{MATSUBARA}{Matsubara, Hayashida and Kano (2018): unified equivalent shear stress amplitude}\label{sheet:MATSUBARA}

\block{Source} G.~Matsubara, A.~Hayashida, D.~Kano, \emph{Int.\ J.\ Fatigue} (2018), accepted
manuscript \cite{matsubara2018}, Eqs.~(7)--(9), (13)--(15); Section~2.2.3.4.1.

\block{Idea and scope} Two parameters describe a multiaxial cycle: an equivalent shear stress
amplitude (crack initiation, Stage~I) and the maximum $S_{max}$ of the normal stress with the
largest range (crack growth, Stage~II, which carries the mean-stress effect). Every state is mapped,
at equal life, onto the fully reversed uniaxial line; there the value is read on the axial S--N
curve. Only axial fatigue data and the true fracture strength are needed, no torsion test.
Proposed for limits and high-cycle lives of steels and aluminium alloys under axial, torsion and
biaxial loading, in and out of phase, with mean stresses and sinusoidal, triangular or trapezoidal
waveforms.

\block{Inputs}
\begin{center}\small
\begin{tabular}{@{}llp{100mm}@{}}
\toprule
Symbol & Unit & Meaning\\\midrule
$s_{ij}$ & MPa & deviatoric stress, $s_{ij}=\sigma_{ij}-\tfrac13\sigma_{kk}\delta_{ij}$ (Eq.~7)\\
$S_1,\dots,S_5$ & MPa & deviatoric components in the basis of Eq.~\eqref{eq:ms-basis}\\
$\sqrt{J_{2,amp}}$ & MPa & equivalent shear stress amplitude (Mamiya), Eq.~\eqref{eq:ms-j2}\\
$S_{max}$ & MPa & maximum over the cycle of the normal stress whose range is the largest over all planes\\
$\sigma_T$ & MPa & true fracture strength\\
$\sigma_w=\fm$ & MPa & fully reversed axial fatigue limit\\
$\sqrt{J_{2,amp,eq}}$ & MPa & unified equivalent shear stress amplitude, Eq.~\eqref{eq:ms-eq}\\
\bottomrule
\end{tabular}
\end{center}

\block{Formulation}
\begin{gather}
S_1=\tfrac{\sqrt3}{2}s_{xx},\quad S_2=\tfrac12(s_{yy}-s_{zz}),\quad S_3=s_{xy},\quad S_4=s_{yz},\quad S_5=s_{zx},
\label{eq:ms-basis}\\
\sqrt{J_{2,amp}}=\sqrt{\sum_{i=1}^{5}\Bigl(\frac{S_{i,\max}-S_{i,\min}}{2}\Bigr)^2} ,
\label{eq:ms-j2}\\
\sqrt{J_{2,amp,eq}} = \frac{\sigma_T\sqrt{J_{2,amp}}}{\sigma_T - S_{max} + \sqrt3\sqrt{J_{2,amp}}} .
\label{eq:ms-eq}
\end{gather}
Eq.~\eqref{eq:ms-eq} is the intersection of the constant-life line through the state, which passes
through $(S_{max},\sqrt{J_2})=(\sigma_T,0)$, with the uniaxial line $\sqrt{J_{2,amp}}/S_{max}=1/\sqrt3$
(Figs.~3--4). On that line $\sqrt{J_{2,amp,eq}}=\sigma_a/\sqrt3$. The fatigue limit (Eq.~9) is
$\sqrt{J_{2,amp,eq}}=\sigma_w/\sqrt3$, i.e.
$\sqrt{J_{2,amp}} + \frac{\sigma_w/\sqrt3}{\sigma_T-\sigma_w}S_{max} = \frac{\sigma_w}{\sqrt3}\frac{\sigma_T}{\sigma_T-\sigma_w}$.

\block{Constants and calibration} $\sigma_T$ from tensile data; when unknown, the paper uses
$\sigma_T=1.84\,S_u$ for $400\le S_u\le1000\MPa$ (its ref.~[23]).

\block{Life} The unified equivalent S--N curve is the reversed axial S--N curve with ordinate
$\sigma_a/\sqrt3$; the life of a state is read at $\sqrt{J_{2,amp,eq}}$. Without axial fatigue data the
paper also estimates this curve from $S_u$ and $\sigma_w$ (Eqs.~16--17).

\block{Repository choices}
\begin{itemize}
\item For one harmonic the half-ranges of $S_i$ are component amplitudes, so $\sqrt{J_{2,amp}}$ is
  the $\sqrt{J_{2,a}}$ of Eq.~\eqref{eq:invariants}.
\item $S_{max}=\sigma_{n,m}+\sigma_{n,a}$ on the plane of the largest $\sigma_{n,a}$, found by a
  search over the sphere of normals; ties go to the larger $S_{max}$ (the paper is silent).
\item $\sigma_T$: workbook constant \texttt{sigma\_T}; otherwise $1.84\,S_u$ only for
  $400\le S_u\le1000\MPa$, else the material is out of scope.
\item Life: the master curve with $\mathrm{DP}=\sqrt{J_{2,amp,eq}}$ fitted on the R~$=-1$ uniaxial
  tests, which is the paper's route; Eqs.~(16)--(17) are not implemented.
\item $S_{max}\ge\sigma_T+\sqrt3\sqrt{J_{2,amp}}$ leaves Eq.~\eqref{eq:ms-eq} undefined and is out of scope.
\end{itemize}

\block{Worked example} Hard steel of Nishihara and Kawamoto ($S_u=680$, $\sigma_w=314$,
$\tm=196\MPa$): $\sigma_T=1.84\cdot680=1251.2\MPa$; fatigue limit $\sigma_w/\sqrt3=181.3\MPa$.
Bending 314: $\sqrt{J_{2,amp}}=181.3$, $S_{max}=314$, $\sqrt{J_{2,amp,eq}}=181.3\MPa$ (onto itself).
Torsion 196: $\sqrt{J_{2,amp}}=S_{max}=196$, $\sqrt{J_{2,amp,eq}}=1251.2\cdot196/(1251.2-196+339.5)=175.8\MPa$,
3\% below the limit. Bending 150.2 with torsion 181.7 at $90^\circ$: $\sqrt{J_{2,amp}}=201.3$,
$S_{max}=199.5$, $\sqrt{J_{2,amp,eq}}=179.9\MPa$.

\block{Verification} The paper reports only error histograms and scatter plots, no per-test value.
Checked: Eq.~\eqref{eq:ms-eq} maps uniaxial states onto $\sigma_a/\sqrt3$ and Eq.~(9) states onto
$\sigma_w/\sqrt3$; Eq.~\eqref{eq:ms-j2} against a time-domain evaluation; $S_{max}$ against a
brute-force search.

% registry: PDS_SWT, PDS_FS
\methodhead{PDS\_SWT, PDS\_FS}{Albinmousa and Al Hussain (2021): polar damage sum}\label{sheet:PDS}

\block{Source} J.~Albinmousa, M.~Al~Hussain, \emph{Forces in Mechanics} 4 (2021) 100025
\cite{albinmousa2021}, Eqs.~(7)--(19); Tables~3 and~5.

\block{Idea and scope} Damage is not confined to the critical plane: every plane perpendicular to
the free surface contributes, and the damage is the area enclosed by the polar curve of a
critical-plane parameter. The critical-plane value is the peak of that curve. Proposed for
strain-controlled low-cycle fatigue (six alloys, 241 tests) under proportional and
non-proportional tension--torsion, with measured stabilised hysteresis loops.

\block{Inputs}
\begin{center}\small
\begin{tabular}{@{}llp{100mm}@{}}
\toprule
Symbol & Unit & Meaning\\\midrule
$\varphi$ & rad & orientation of a plane perpendicular to the surface ($x$--$y$ plane)\\
$\sigma_\varphi$, $\varepsilon_\varphi$, $\gamma_\varphi$ & MPa, --, -- & normal stress, normal strain, engineering shear strain on plane $\varphi$\\
$\sigma_{\max}(\varphi)$, $\sigma_{n,\max}(\varphi)$ & MPa & maximum over the cycle of $\sigma_\varphi$\\
$\Delta\varepsilon_{\max}(\varphi)$, $\Delta\gamma_{\max}(\varphi)$ & -- & ranges of $\varepsilon_\varphi$ and $\gamma_\varphi$\\
$\sigma_{yt}$ & MPa & tensile yield strength\\
$k_\varphi$ & -- & Fatemi--Socie constant of the polar sum\\
$\nu^{el}$, $\nu^{pl}$ & -- & elastic and plastic Poisson's ratios\\
$D(\varphi)$, $D^s$ & & plane damage and polar damage sum\\
$\alpha$, $\beta$ & & constants of the life law\\
\bottomrule
\end{tabular}
\end{center}

\block{Formulation} With $\sigma_x=0$, the plane transformations are (Eqs.~7a--7e)
\begin{gather}
\sigma_\varphi = \tfrac{\sigma_y}{2}+\tfrac{\sigma_y}{2}\cos2\varphi-\tau_{xy}\sin2\varphi ,\qquad
\varepsilon_\varphi = \tfrac{\varepsilon_x+\varepsilon_y}{2}-\tfrac{\varepsilon_x-\varepsilon_y}{2}\cos2\varphi-\tfrac{\gamma_{xy}}{2}\sin2\varphi ,\\
\tfrac{\gamma_\varphi}{2} = \tfrac{\gamma_{xy}}{2}\cos2\varphi-\tfrac{\varepsilon_x-\varepsilon_y}{2}\sin2\varphi ,\qquad
\varepsilon_x = -(\nu^{el}\varepsilon_y^{el}+\nu^{pl}\varepsilon_y^{pl}) .
\end{gather}
The plane damage is the SWT (Eq.~8) or the Fatemi--Socie (Eq.~9) parameter,
\begin{equation}
D_{SWT}(\varphi) = \sigma_{\max}(\varphi)\,\frac{\Delta\varepsilon_{\max}(\varphi)}{2} ,\qquad
D_{FS}(\varphi) = \frac{\Delta\gamma_{\max}(\varphi)}{2}\Bigl(1+k_\varphi\frac{\sigma_{n,\max}(\varphi)}{\sigma_{yt}}\Bigr) ,
\end{equation}
and the polar damage sum is the enclosed area (Eqs.~10--14),
\begin{equation}
D^s = \frac12\int_0^{2\pi} D(\varphi)^2\,d\varphi .
\label{eq:pds}
\end{equation}

\block{Constants and calibration} $k_\varphi$ is varied until the $D^s$--life points of the axial
and shear tests fall on one band (Fig.~4); $\alpha$, $\beta$ by a power fit. Table~3 reports
$k_\varphi$ from 0.1 to 5.0.

\block{Life} A power law on the polar sum, $D^s=\alpha_n N_f^{\beta_n}$ (SWT) and
$D^s=\alpha_s N_f^{\beta_s}$ (FS), Eqs.~(17)--(19), fitted on axial and shear tests together.

\block{Repository choices}
\begin{itemize}
\item \emph{Elastic version} (maintainer decision): the strains are the linear-elastic strains of
  the stress (Section~\ref{sec:notation}), not measured loops; this is a high-cycle reading of a
  low-cycle method.
\item For one harmonic, on $\n=(\cos\psi,\sin\psi,0)$ with tangent $\mathbf t$:
  $\sigma=\n^{\mathsf T}\Tsig\n$, $\varepsilon=\n^{\mathsf T}\Teps\n$,
  $\gamma=2\,\mathbf t^{\mathsf T}\Teps\n$; maxima and half ranges from Eq.~\eqref{eq:scalar}.
  The integral uses 720 equally spaced planes.
\item $D(\varphi)$ is replaced by $\max(D(\varphi),0)$ before squaring.
\item $k_\varphi$ (\texttt{k\_PDS}) is fitted like $k_{FS}$: minimum life scatter of the R~$=-1$
  tension and torsion failures, $0\le k_\varphi\le10$; on fatigue-limit sets it is set so that both
  limits give the same $D^s$. A value at the search bound is not identifiable.
\item Life: the master curve on $D^s$, fitted on the uniaxial tests (the paper fits axial and shear
  together).
\item Plane stress on the surface only. An in-phase equibiaxial state has no in-plane shear, so
  $D_{FS}=0$ and its load factor is undefined.
\end{itemize}

\block{Worked example} Steel, $E=206\,000\MPa$, $\nu=0.3$, $\sigma_{yt}=727\MPa$, $k_\varphi=1$;
paths with the same von Mises amplitude $400\MPa$:
\begin{center}\small
\begin{tabular}{@{}lccc@{}}
\toprule
Path & $D^s_{SWT}$ (MPa$^2$) & peak $D_{SWT}$ (MPa) & $D^s_{FS}$\\\midrule
tension 400 & 0.4778 & 0.777 & $1.646\cdot10^{-5}$\\
torsion 231 & 0.1335 & 0.337 & $1.728\cdot10^{-5}$\\
tension 283 + torsion 163 at $90^\circ$ & 0.2124 & 0.389 & $1.862\cdot10^{-5}$\\
\bottomrule
\end{tabular}
\end{center}
For tension the peak is $\sigma_a^2/E=0.777\MPa$, the critical-plane SWT value.

\block{Verification} The paper's results rest on measured LCF loops; no per-test value is printed.
Checked: Eqs.~(7a), (7c), (7d); the uniaxial closed form and its peak; convergence of the sum to
Eq.~\eqref{eq:pds}.

% registry: WANG2018
\methodhead{WANG2018}{Wang, Shang, Sun and Chen (2018): weighted critical plane with mean shear stress}\label{sheet:WANG2018}

\block{Source} X.-W.~Wang, D.-G.~Shang, Y.-J.~Sun, H.~Chen, \emph{Theor.\ Appl.\ Fract.\ Mech.}
(2018), accepted manuscript \cite{wang2018}, Eqs.~(1)--(20); Tables~1--5; Section~5.1.

\block{Idea and scope} The critical plane is the weighted time average of the planes of maximum
shear stress. The weights grow with the instantaneous maximum shear stress. On that plane, the
shear and normal stress amplitudes are combined in von Mises form. Each is corrected by its own
mean stress: the shear mean through its absolute value, the normal mean with its sign. The life is
read on the fully reversed torsion S--N curve. Proposed for high-cycle lives of 7075-T651
thin-walled tubes under tension--torsion: constant amplitude, in and out of phase, with mean
stresses, and variable amplitude.

\block{Inputs}
\begin{center}\small
\begin{tabular}{@{}llp{100mm}@{}}
\toprule
Symbol & Unit & Meaning\\\midrule
$\sigma(t)$, $\tau(t)$ & MPa & axial normal stress and shear stress of the tube ($\sigma_{yy}$, $\tau_{xy}$ in the repository frame)\\
$\varphi$ & rad & angle between the normal of a plane perpendicular to the surface and the axis\\
$\sigma_\varphi$, $\tau_\varphi$ & MPa & normal and shear stress on plane $\varphi$, Eq.~\eqref{eq:w18-plane}\\
$\tau_{\max}(t_k)$ & MPa & maximum shear stress over the surface planes at instant $t_k$\\
$\varphi_{45}(t_k)$ & rad & orientation of the maximum-shear plane at $t_k$, in $[-45^\circ,45^\circ]$\\
$W(t_k)$ & -- & weight of instant $t_k$, Eq.~\eqref{eq:w18-weight}\\
$\hat\varphi$, $\varphi_c$ & rad & weighted maximum-shear plane and critical plane\\
$\tau_a,\tau_m$; $\sigma_a,\sigma_m$ & MPa & amplitude and mean of $\tau_\varphi$ and $\sigma_\varphi$ on $\varphi_c$\\
$\tau_{-1}=\tm$, $\sigma_{-1}=\fm$ & MPa & fully reversed torsion and axial limits\\
$\sigma_0$, $\tau_0$ & MPa & pulsating ($R=0$) axial and torsion limits, maximum stress\\
$k_\tau$, $k_\sigma$ & -- & mean shear and mean normal stress sensitivities, Eq.~\eqref{eq:w18-k}\\
$\tau_{eq}$ & MPa & equivalent shear stress, Eq.~\eqref{eq:w18}\\
\bottomrule
\end{tabular}
\end{center}

\block{Formulation} On the planes perpendicular to the surface (Eq.~2),
\begin{equation}
\sigma_\varphi = \tfrac{\sigma}{2}+\tfrac{\sigma}{2}\cos2\varphi+\tau\sin2\varphi ,\qquad
\tau_\varphi = \tfrac{\sigma}{2}\sin2\varphi-\tau\cos2\varphi .
\label{eq:w18-plane}
\end{equation}
At each instant $t_k$ of one cycle (Eqs.~3--7),
\begin{gather}
\tau_{\max}(t_k)=\sqrt{\bigl(\sigma(t_k)/2\bigr)^2+\tau(t_k)^2},\qquad
\varphi_{45}(t_k)=\tfrac12\arctan\!\left[-\frac{\sigma(t_k)}{2\tau(t_k)}\right],\\
W(t_k)=\frac{\tau_{\max}(t_k)-\min_k\tau_{\max}}{\max_k\tau_{\max}-\min_k\tau_{\max}},\qquad
\lvert\hat\varphi\rvert=\frac{\sum_k\lvert\varphi_{45}(t_k)\rvert\,W(t_k)}{\sum_k W(t_k)} ,
\label{eq:w18-weight}
\end{gather}
and $\hat\varphi$ takes the sign of $\sum_k\varphi_{45}(t_k)W(t_k)$. The shear stress has the same
magnitude on $\hat\varphi$ and $\hat\varphi+90^\circ$, so the candidate with the larger damage is
$\varphi_c$. When $\tau_{\max}$ is constant (Eqs.~8--9), the plane of the largest normal stress is
critical. On $\varphi_c$ (Eqs.~11, 15, 18),
\begin{gather}
\tau_{eq}=\sqrt{\bigl[\tau_a+k_\tau\lvert\tau_m\rvert\bigr]^2+\tfrac13\bigl[\sigma_a+k_\sigma\sigma_m\bigr]^2}\le\tau_{-1},
\label{eq:w18}\\
k_\tau=\frac{2\tau_{-1}}{\tau_0}-1,\qquad
k_\sigma=\sqrt{\frac{48\tau_{-1}^2}{\sigma_0^2}-\frac{12\tau_{-1}^2}{\tau_0^2}}-1 ,
\label{eq:w18-k}
\end{gather}
obtained by imposing reversed torsion ($\tau_a=\tau_{-1}$), pulsating tension
($\tau_a=\tau_m=\sigma_a=\sigma_m=\sigma_0/4$ on the $45^\circ$ plane) and pulsating torsion
($\tau_a=\tau_m=\tau_0/2$), Eqs.~(12)--(14). Eq.~(17) requires $k_\tau>0$ and $k_\sigma>0$.
Without the pulsating limits (Eq.~19),
\begin{equation}
\sigma_0=\sqrt2\,\sigma_{-1},\qquad \tau_0=\frac{4\tau_{-1}}{2\sigma_{-1}/\sigma_0+1},
\end{equation}
which gives $k_\tau=(\sqrt2-1)/2$ (Eq.~20). Reversed tension gives $\tau_{eq}=\sigma_a/\sqrt3$: the
criterion implies $\sigma_{-1}=\sqrt3\,\tau_{-1}$.

\block{Constants and calibration} $k_\tau$, $k_\sigma$ from $\tau_{-1}$, $\sigma_0$, $\tau_0$, or
from $\sigma_{-1}$, $\tau_{-1}$ through Eq.~(19). The paper takes $\sigma_{-1}$ and $\tau_{-1}$
from its S--N curves at $N_0=10^6$, $2\cdot10^6$ and $10^7$, and compares the three choices.

\block{Life} $\tau_{eq}$ on the fully reversed torsion S--N curve (Section~4.2). Variable
amplitude: Wang--Brown reversals on the two candidate planes, Miner's rule, and the larger damage
(Eqs.~21--24).

\block{Repository choices}
\begin{itemize}
\item Scope: tension--torsion only (maintainer decision). Any other stress component is out of
  scope.
\item One cycle sampled at 3600 instants for Eqs.~\eqref{eq:w18-weight}. On a plane, the
  amplitude and mean are exact single-harmonic values.
\item The two candidates are compared through $\tau_{eq}$: there is one S--N curve, so the larger
  $\tau_{eq}$ is the larger damage. A zero signed sum gives $\hat\varphi\ge0$.
\item Constant $\tau_{\max}$: the plane of the largest $\max_t\sigma_\varphi(t)$ is found over
  3600 planes.
\item $\sigma_{-1},\tau_{-1}$ are the workbook \texttt{Sigm1}, \texttt{Taum1}. $\sigma_0,\tau_0$
  are \texttt{Sig0}, \texttt{Tau0}, otherwise Eq.~(19). A violation of Eq.~(17) is out of material
  scope.
\item Life: the master curve $\tau_{eq}=AN^b$ is fitted on the $R=-1$ pure torsion failures,
  which is the paper's torsion S--N curve. The variable-amplitude procedure is not implemented.
\end{itemize}

\block{Worked example} 7075-T651, Table~2 curves at $N_0=10^6$: $\sigma_{-1}=222.0$,
$\tau_{-1}=127.1\MPa$, so $k_\tau=0.2071$ and $k_\sigma=0.8698$. Table~5 load
$\sigma=163.02+144.57\sin\omega t$, $\tau=94.12+83.47\sin(\omega t-\delta)$ (MPa):
\begin{center}\small
\begin{tabular}{@{}lccc@{}}
\toprule
$\delta$ & $\varphi_c$ & $\tau_{eq}$ (MPa) & $N$ from $\tau=688.85N^{-0.12233}$\\\midrule
$0^\circ$ & $69.6^\circ$ & 159.3 & 157\,636\\
$30^\circ$ & $-20.4^\circ$ & 160.0 & 152\,280\\
$90^\circ$ & $-20.1^\circ$ & 160.4 & 149\,120\\
\bottomrule
\end{tabular}
\end{center}
Experimental lives: 53\,058--65\,046, 64\,231--82\,445 and 70\,333--94\,718 cycles.

\block{Verification} Section~5.1: the estimated crack trace, $90^\circ-\lvert\varphi_c\rvert$, is
$70^\circ$ for every $30^\circ$ and $90^\circ$ out-of-phase test of Table~5, as in the paper.
Eqs.~(12)--(14) give $\tau_{eq}=\tau_{-1}$ exactly; Eq.~(20) equals Eq.~(15) with Eq.~(19). The
paper prints the predicted lives only in figures.


\bibliographystyle{unsrt}
\bibliography{references}
\end{document}
